\begin{abstract}
Zero-shot object-goal navigation agents driven by vision--language
models lose much of their success rate once the camera input is
corrupted by low light, haze, or motion, and a common remedy is to
restore the image before it reaches the policy. Through a controlled
study on InstructNav in Habitat--Matterport 3D ObjectNav, with a
same-seed paired protocol at $N{=}300$ and a blind per-step reliability
estimator, we show that this remedy targets the wrong layer.
Degradation does not stop the agent from acting; it corrupts the
quality of an irreversible commitment: the agent still reaches within
one metre of the goal in $61\%$ of episodes yet succeeds in only
$37\%$, and restoration that visibly repairs the image moves success
by a non-significant $+0.04$. A sweep over six causal families of
test-time interventions---pixel restoration with a per-frame appearance
upper bound, reliability-gated reweighting, commit-timing revocation,
forced commitment, temporal/multi-view aggregation, and a depth-based
multimodal probe---finds that every deployable family lands at
${\sim}0$ or significantly hurts success, evidence for a commitment
error that none of the probed test-time channels recovers.
The fingerprint survives a cross-corruption replication
(low-light, motion-blur, fog, Gaussian noise) on the same
InstructNav+Qwen2-VL-7B pipeline. These
results point the field away from test-time patches on frozen,
clean-trained agents, toward degradation-aware training and learned
multi-modal sensing. We release the benchmark, the paired protocol,
and all of the negative results.
\end{abstract}